\subsection{Sidebar: SymPy project}
%wrap figure
%\begin{wrapfigure}{r}{.2\textwidth}
\begin{figure}[h]
\vspace{-0pt}
\centering
\includegraphics[width=.18\textwidth]{images/logo}
\vspace{-0pt}
\label{fig:}
\vspace{00pt}
\end{figure}
%\end{wrapfigure}

\begin{description}
\item {\bf What is SymPy?}\\
SymPy is a Python library for symbolic mathematics. It aims to become a
full-featured computer algebra system (CAS) while keeping the code as simple as
possible in order to be comprehensible and easily extensible. SymPy is written
entirely in Python and does not require any external libraries.
\item{\bf Where do I find it?}\\
\begin{tabular}{ll}
{\bf Website}: & http://sympy.org \\
{\bf Mailing List}: & http://groups.google.com/group/sympy \\
{\bf IRC}: & irc://irc.freenode.net/sympy \\
{\bf Source}: & http://github.com/sympy/sympy
\end{tabular}
\item {\bf Who develops SymPy?}\\
SymPy has been developed since 2006 as an open source project under the bazaar
model with over 150 contributors.  Surprisingly, it is not developed by any
particular university group or company but has become a competitive alternative
to many CAS projects that are.  Google Summer of Code has funded 25
student contributors over the summers since 2007.
\item{\bf Who uses SymPy?}\\
Because SymPy is open source and pure Python, it is easy to integrate its
functionality into other research projects. Several projects exist which make
use of SymPy in different disciplines. Examples include Sage, PyDy, Ignition,
several symbolic math GUI's, and many more.
\item{\bf How can you contribute?}\\
While no project is perfect, SymPy could be a little faster, better documented
and allow for much easier embedding in other projects.  New wanted features
include new solvers, new symbolic integration algorithms, and better domain
support.  For those interested in helping see the development section of the
website.
\end{description}
